(1) 在圆轨道旋转系中，Coriolis项$-2\boldsymbol n\times\dot{\boldsymbol r}$与离心项$-\boldsymbol n\times(\boldsymbol n\times\boldsymbol r)$必须同时保留。把引力梯度$-GM[\boldsymbol r-3(\hat r\cdot\boldsymbol r)\hat r]/r_0^3$与原圆轨道平衡相减，得到$\ddot x-2n\dot y-3n^2x=0$、$\ddot y+2n\dot x=0$、$\ddot z+n^2z=0$，这就是HCW线性模型。
(2) 对平面方程作常系数求解，可写$x(t)=A\cos nt+B\sin nt+C$，$y(t)=-2A\sin nt+2B\cos nt-\frac32nCt+D$。独立法向运动为$z(t)=E\cos nt+F\sin nt$。逐项代回三条微分方程能验证形式完整：平面含两振荡常数、一个线性漂移常数和一个位置常数。
(3) 从$t=0$初值，$x_0=A+C$且$\dot y_0=-2nA-\frac32nC$。消去$A$得$C=4x_0+2\dot y_0/n$，因此消除无界顺轨漂移必须$C=0$，即$\dot y_0=-2nx_0$。这一条件并不保证两颗卫星相互静止，只确保相对轨迹在该线性近似中不因半长轴偏差持续漂移。
(4) 将HCW方程分别乘$\dot x,\dot y,\dot z$后相加，Coriolis做功抵消，得到$dJ/dt=0$，其中$J=\frac12(\dot x^2+\dot y^2+\dot z^2)-\frac32n^2x^2+\frac12n^2z^2$。该Jacobi型量不是正定能量，允许漂移模；当相对距离不再小于$r_0$或轨道明显非圆时，必须使用非线性三体或椭圆参考轨道方程。
